\documentclass[11pt,a4paper]{article}
\usepackage[margin=22mm]{geometry}
\usepackage{amsmath,amssymb,booktabs,array}
\usepackage[T1]{fontenc}
\usepackage{lmodern}
\usepackage{xcolor}
\usepackage{hyperref}
\hypersetup{colorlinks=true,linkcolor=blue!50!black,urlcolor=blue!50!black}
\setlength{\parindent}{0pt}
\setlength{\parskip}{6pt}
\renewcommand{\arraystretch}{1.18}
\newcommand{\E}{\mathbb{E}}
\newcommand{\sourcefile}[1]{\texttt{\detokenize{#1}}}
\title{Reconstruction of a Jump-Diffusion Option Pricing Study\\[4pt]
\large Implemented Components and Verified Results}
\author{Project Replication Report}
\date{4 October 2026}

\begin{document}
\maketitle
\vspace{-10mm}
\begin{abstract}
This report records the parts recreated from the study by
Volk-Makarewicz, Borovkova, and Heidergott~\cite{paper}: the experimental
design, model and payoff constructions, local gradient calculations,
and figure definitions. It also presents independent checks of the
implementation and one close numerical comparison with the paper.
The scope is an implemented and checked research baseline; complete
numerical agreement with the published graphs has not been established.
Unmatched curves are excluded from this report.
\end{abstract}

\section{Purpose and recreated experiment design}
The study assesses how adding jumps changes option prices and whether
gradient estimates can detect that change. The reconstruction compares
a diffusion model, denoted by $P$, with its jump extension, denoted by $Q$.
Its intended use is to provide a baseline before investigating alternative
approaches.

\begin{table}[h]
\centering\small
\begin{tabular}{@{}p{39mm}p{38mm}p{46mm}r@{}}
\toprule
Baseline & Jump extension & Option & Dates\\
\midrule
Geometric Brownian motion & Merton & European call & 1\\
Vasicek & Mean-reverting jump diffusion & European put & 1\\
Vasicek & Mean-reverting jump diffusion & Arithmetic Asian call & 4\\
Vasicek & Mean-reverting jump diffusion & Arithmetic Asian put & 12\\
\bottomrule
\end{tabular}
\caption{Implemented option cases. The twelve-date case follows the
paper's payoff formula and panel heading; its caption is inconsistent.}
\end{table}

The reconstructed experiment uses maturity $T=1$, rate $r=0$, and jump
intensities $\lambda\in\{0.01,0.08,0.70\}$. The Merton case uses
$S_0=K=40$; the mean-reverting cases use $S_0=40$ and $K=35$.
There are 500 outer repetitions per Merton setting and 300 per
mean-reverting setting. These are independent simulation runs; paths
within a run use variance-reduction sampling.

For the ratio experiments, the Merton case uses 16,384 paths per estimate
and the mean-reverting cases use 256. The path-count experiment uses
$64,128,\ldots,16{,}384$ paths, with fixed ratios $0.5$ and $0.33$,
respectively. The test significance level is $\alpha=0.10$.

\textbf{Presentation takeaway.} The option cases, experiment grids, path
counts, and outer repetition counts have been recreated. Matching this
design alone does not establish matching numerical curves.

\newpage
\section{Implemented model and payoff constructions}
\subsection{GBM and Merton}
The diffusion baseline is simulated with exact log-price increments.
The implemented Merton extension has the terminal construction
\begin{align}
S_t^{P}&=S_0\exp\!\left[(r-\tfrac12\sigma^2)t+\sigma W_t\right],\\
S_t^{Q}&=S_0\exp\!\left[(r-\tfrac12\sigma^2-\lambda\eta)t
                  +\sigma W_t+\sum_{k=1}^{N_t}Y_k\right],
\end{align}
where $N_t$ has Poisson intensity $\lambda$,
$Y_k\sim\mathcal N(\mu_J,s_J^2)$, and
$\eta=\E[e^{Y_k}-1]=e^{\mu_J+s_J^2/2}-1$.
The configured mark parameters are $\mu_J=-3.3\times10^{-5}$ and
$s_J=0.03347$. For the positive-intensity ratio sweep,
$\sigma=\lambda/\rho$, with
\[
\rho\in\{0.25,0.40,0.50,0.63,0.79,1.00,1.59,2.00\}.
\]
The mapping of the paper's $\mu_J$ to the mean log jump is an explicit
interpretation because that symbol is also used for a drift quantity.

\subsection{Vasicek and mean-reverting jump diffusion}
With the configured diffusion risk adjustment set to zero, the
implemented mean-reverting models can be written as
\begin{align}
dX_t^{P}&=\kappa(\theta-X_t^{P})\,dt+\sigma\,dW_t,\\
dX_t^{Q}&=\bigl[\kappa(\theta-X_t^{Q})-\lambda\E[J]\bigr]dt
             +\sigma\,dW_t+dL_t,
\end{align}
where $L_t=\sum_{k=1}^{N_t}J_k$. Exact Ornstein--Uhlenbeck transitions
are used between observation dates. The configured values are
$\kappa=0.24$, $\theta=0.63$, and $\sigma=5.16$.
Positive and negative marks each have probability $1/2$, with exponential
mean magnitudes $\beta_+$ and $\beta_-$. Thus
\[
\E[J]=\tfrac12(\beta_+-\beta_-),\qquad
\E[|J|]=\tfrac12(\beta_++\beta_-).
\]
The baseline magnitudes are $3.48$ and $5.54$. For the plotted ratio
$\rho=\E[|J|]/\sigma$, the revised mapping preserves the signed mean $m_J=-1.03$ over
$\{0.20,0.25,0.29,0.33,0.40,0.50,0.67,1.00\}$.
With $a=\rho\sigma$ and positive-jump probability $p$, it sets
\[
\beta_+=\frac{a+m_J}{2p},\qquad
\beta_-=\frac{a-m_J}{2(1-p)}.
\]
This reconciles the figure axis with the statement that the signed mean
stays fixed. It requires $a>|m_J|$, satisfied by every configured ratio.
The authors' executed mapping remains unverified.

\subsection{Discounted payoffs}
At dates $t_j=jT/m$, define $A_m=m^{-1}\sum_{j=1}^m X_{t_j}$.
European payoffs use the terminal asset value; arithmetic Asian payoffs
use $A_m$. Calls and puts are evaluated as
\[
H_{\mathrm{call}}(x)=e^{-rT}(x-K)^+,
\qquad H_{\mathrm{put}}(x)=e^{-rT}(K-x)^+.
\]
Merton log states are exponentiated before evaluating payoffs.
Mean-reverting states are used directly. Discounting is applied once to
prices and gradient contributions.

\newpage
\section{Gradient construction and implementation corrections}
For a simulated jump-time partition, let $B$ be the baseline path state
and let $D_i$ be the state change caused by replacing interval $i$ with
its jump-model counterpart. Each $D_i$ includes the local drift
adjustment and any jump contribution. Using the discounted payoff $H$,
the implemented exhaustive constructions are
\begin{align}
G_1&=\sum_i\bigl[H(B+D_i)-H(B)\bigr],\\
G_2&=G_1+\sum_{i<j}\bigl[H(B+D_i+D_j)
             -H(B+D_i)-H(B+D_j)+H(B)\bigr].
\end{align}
The fully replaced payoff is $H(B+\sum_iD_i)$. GradOne uses the singleton
sum; GradTwo adds unordered pair contributions. They represent the
first- and second-order local expansion constructions. They are not
assumed to equal the full price difference in every setting.

The common-random-number comparison uses paired differences
$H_Q-H_P$. Baseline and jump-model payoffs share the underlying
randomness. Brownian increments and jump times use antithetics;
jump occurrence is stratified. The exhaustive estimator is the current
baseline. Enumeration over interval choices at 72 method settings verifies
that the randomized alternative has the same conditional expectation
(maximum discrepancy below $3\times10^{-13}$). Its added randomization
variance is recorded separately. The authors' exact randomization remains
unverified.

\subsection{Recreated figure definitions}
\begin{table}[h]
\centering\small
\begin{tabular}{@{}cp{30mm}p{39mm}p{49mm}@{}}
\toprule
Figure & Horizontal axis & Upper panels & Lower panels\\
\midrule
1 & Model ratio & Median estimated power & Relative difference of median prices\\
2 & Paths per estimate & Empirical rejection frequency & Median estimated power\\
3 & Same ratio grid as Figure 1 & Mean estimated power & Relative difference of mean prices\\
\bottomrule
\end{tabular}
\caption{Recreated plot structure. Numerical power curves are outside
the verified scope of this report.}
\end{table}

Figures 1 and 3 now share the ratio-sweep observations, with different
aggregation rules. Figure 3 previously used the path-count experiment.
Price observations are counted once per replicate, rather than being
duplicated across the three method records.

\subsection{Completed code corrections}
\begin{itemize}
\item Made the MRJD mapping explicit and preserved its signed mean while varying the absolute mean.
\item Corrected Figure 3's data source and implemented the exhaustive
singleton and pairwise gradient sums.
\item Applied consistent discounting and retained the configured
diffusion volatility when Merton jump intensity is zero.
\item Vectorized jump-arrival and interval-effect calculations;
retained exact diffusion transitions.
\item Reduced verbose comments, formatted the source, and strengthened
configuration and payoff checks.
\item Preserved the legacy simulation and previous outputs, and recorded
seeds, environment details, and file hashes for the corrected run. DEBUG outputs are isolated from
the full experiment.
\end{itemize}

\newpage
\section{Checks and the close numerical result}
\subsection{Independent implementation checks}
The recorded analytical checks use 100,000 paths per check. They compare
the simulation with closed-form prices or moments, rather than with
another execution of the same path construction.
\begin{table}[h]
\centering\small
\begin{tabular}{@{}p{73mm}rr@{}}
\toprule
Quantity & Analytical & Simulation\\
\midrule
GBM call price (Black--Scholes) & 0.270958 & 0.270023\\
Merton call price (Poisson mixture) & 2.554915 & 2.552704\\
Vasicek terminal mean & 31.599539 & 31.599539\\
Vasicek terminal variance & 21.146085 & 21.087735\\
Compensated MRJD terminal mean & 31.599539 & 31.600677\\
MRJD terminal variance & 21.533752 & 21.466882\\
\bottomrule
\end{tabular}
\caption{All six checks passed the implementation's recorded criteria.
These checks use their own benchmark settings, not necessarily the
close paper endpoint below.}
\end{table}

The test suite records \textbf{35 passing tests}, including independently
enumerated local mixture expansions, compensated transitions, payoff
checks, and independent option-price benchmarks. Ruff checks passed.
Characteristic-function quadrature also checks MRJD European and
arithmetic Asian prices under the implemented assumptions. Across the
ratio-sweep settings, maximum absolute differences between pooled
simulation mean prices and quadrature prices were about $0.01081$ for
$P$ and $0.01355$ for $Q$.

These results support internal correctness under the chosen model
interpretations. They do not establish agreement with all published
curves. In particular, antithetic and stratified paths are dependent;
the recorded path-level standard-error convention is not independently
certified by these checks.

\subsection{Close Merton price-difference endpoint}
For Figure 1(a), at $\lambda=0.01$ and $\lambda/\sigma=2$, the relative
difference of median price estimates is
\[
R_{\mathrm{med}}=
\frac{|\operatorname{median}(\widehat V_Q)
       -\operatorname{median}(\widehat V_P)|}
     {|\operatorname{median}(\widehat V_P)|}.
\]
\begin{center}
\begin{tabular}{@{}lr@{}}
\toprule
Approximate visual reading from the paper & $0.058$\\
Recreated result & $0.057702$\\
Absolute difference & $0.000298$\\
\bottomrule
\end{tabular}
\end{center}
The discrepancy is approximately $0.51\%$ of the visual paper reading.
The source value is approximate, so this is evidence of close agreement
at one endpoint, not an exact numerical reproduction of the full graph.

\section{Scope for presentation use}
The completed run contains 214,200 method records and 612 summary
records, using the published outer repetition counts. Its recorded
simulation runtime is approximately 8.1 minutes.

The defensible presentation claims are: the experimental design and
figure definitions were recreated; the model, payoff, and gradient
implementation was corrected and checked; and one Merton price
endpoint closely agrees with the supplied paper. MRJD published price
curves and the power curves remain unverified and are omitted here.
A separate resolution register records the paper's conflicting ratio,
contract captions, derivative coefficients, no-jump statement, and inference
conventions, with evidence and the selected implementation choices.
The full study of alternative specifications contains 120 parameter
settings and 480 independent expected-price comparisons.

\small
\begin{thebibliography}{1}
\bibitem{paper} W. Volk-Makarewicz, S. Borovkova, and B. Heidergott.
Assessing the impact of jumps in an option pricing model: A gradient
estimation approach. \emph{European Journal of Operational Research},
298:740--751, 2022.
\href{https://doi.org/10.1016/j.ejor.2021.07.015}{doi:10.1016/j.ejor.2021.07.015}.
\end{thebibliography}
\textbf{Local evidence:}
\sourcefile{config/experiments.json};
\sourcefile{outputs/logs/run_metadata.json};
\sourcefile{outputs/validation/figure_comparison.csv};
\sourcefile{outputs/validation/simulation_vs_analytic_prices.csv};
\sourcefile{reproduction/reproducibility_manifest.json}.
\end{document}
